% GENERATED FILE. Edit canonical lessons, then run npm run generate:books.
% canonical-source-hash: fnv1a64:bc2ae80c99d4f654
% canonical-lessons: GU-C404-phutpath, GU-C404-kedi, GU-C404-sakdu, GU-C404-cadhan, GU-C404-tapvu-leap

\chapter{Pavements and Paths}
\label{ch:gu-pavements-and-paths}
\hlchaptermodality{404}

\begin{chapteropening}
\textbf{By the end of this chapter:} \emph{I can describe a path for walking.}

Complete the last lesson of chapter 404: I can describe a path for walking.
\end{chapteropening}

\section[\texorpdfstring{\gu{ફૂટપાથ} (phūṭpāth)}{ફૂટપાથ (phūṭpāth)}]{\gu{ફૂટપાથ} (phūṭpāth) — a pavement; a footpath (beside a road)}

\label{lesson:GU-C404-phutpath}



\begin{quote}
Before the new one: say the Gujarati for to branch off, then the Gujarati for to crash (into).
\end{quote}

\subsection*{You'll want to know: \gu{ફૂટપાથ}}
\textbf{\gu{ફૂટપાથ}} — \emph{phūṭpāth} — \textquotedblleft{}a pavement; a footpath (beside a road)\textquotedblright{}.

\begin{quote}
\gu{બાળકોએ} \gu{ફૂટપાથ} \gu{પર} \gu{ચાલવું} \gu{જોઈએ}.
\end{quote}

\begin{quote}
\emph{bāḷkoe phūṭpāth par cālvũ joīe.}
\end{quote}

\begin{quote}
Children should walk on the pavement.
\end{quote}

A feminine noun.

\begin{cousinweb}
\textbf{\gu{ફૂટપાથ}} (\emph{phūṭpāth}) is borrowed from English \emph{footpath}, the raised walk beside a city road.
\end{cousinweb}

\begin{grammarlens}[title={What you've built}]
One more word for getting around.
\end{grammarlens}

\subsection*{Your turn}
\begin{itemize}
\raggedright
  \item \emph{Say it:} \emph{phūṭpāth}
  \item \emph{Say it:} \emph{phūṭpāth}, once more
  \item \emph{Say it:} \emph{bāḷkoe phūṭpāth par cālvũ joīe.} — Children should walk on the pavement.
  \item \emph{Recall:} say the Gujarati for to branch off, then the Gujarati for to crash (into), then say \emph{phūṭpāth} again
\end{itemize}

\begin{tcolorbox}[breakable,colback=teal!4,colframe=teal!35!black,title={Before you move on}]
What does \gu{ફૂટપાથ} mean? (\textquotedblleft{}A pavement; a footpath (beside a road)\textquotedblright{}.) Say it once more.
\end{tcolorbox}
\section[\texorpdfstring{\gu{કેડી} (keḍī)}{કેડી (keḍī)}]{\gu{કેડી} (keḍī) — a footpath (across fields); a trail}

\label{lesson:GU-C404-kedi}



\begin{quote}
Before the new one: say the Gujarati for to crash (into), then the Gujarati for a pavement.
\end{quote}

\subsection*{You'll want to know: \gu{કેડી}}
\textbf{\gu{કેડી}} — \emph{keḍī} — \textquotedblleft{}a footpath (across fields); a trail\textquotedblright{}.

\begin{quote}
\gu{ખેતરની} \gu{વચ્ચે} \gu{એક} \gu{નાની} \gu{કેડી} \gu{છે}.
\end{quote}

\begin{quote}
\emph{khetarnī vacce ek nānī keḍī che.}
\end{quote}

\begin{quote}
There is a little trail through the middle of the field.
\end{quote}

A feminine noun.

\begin{cousinweb}
\textbf{\gu{કેડી}} (\emph{keḍī}) No dictionary we use records an older form, so its history is not settled.
\end{cousinweb}

\begin{grammarlens}[title={What you've built}]
One more word for getting around.
\end{grammarlens}

\subsection*{Your turn}
\begin{itemize}
\raggedright
  \item \emph{Say it:} \emph{keḍī}
  \item \emph{Say it:} \emph{keḍī}, once more
  \item \emph{Say it:} \emph{khetarnī vacce ek nānī keḍī che.} — There is a little trail through the middle of the field.
  \item \emph{Recall:} say the Gujarati for to crash (into), then the Gujarati for a pavement, then say \emph{keḍī} again
\end{itemize}

\begin{tcolorbox}[breakable,colback=teal!4,colframe=teal!35!black,title={Before you move on}]
What does \gu{કેડી} mean? (\textquotedblleft{}A footpath (across fields); a trail\textquotedblright{}.) Say it once more.
\end{tcolorbox}
\section[\texorpdfstring{\gu{સાંકડું} (sā̃kḍũ)}{સાંકડું (sā̃kḍũ)}]{\gu{સાંકડું} (sā̃kḍũ) — narrow}

\label{lesson:GU-C404-sakdu}



\begin{quote}
Before the new one: say the Gujarati for a pavement, then the Gujarati for a trail.
\end{quote}

\subsection*{You'll want to know: \gu{સાંકડું}}
\textbf{\gu{સાંકડું}} — \emph{sā̃kḍũ} — \textquotedblleft{}narrow\textquotedblright{}.

\begin{quote}
\gu{આ} \gu{ગલી} \gu{એટલી} \gu{સાંકડી} \gu{છે} \gu{કે} \gu{ગાડી} \gu{અંદર} \gu{નથી} \gu{જતી}.
\end{quote}

\begin{quote}
\emph{ā galī eṭlī sā̃kḍī che ke gāḍī andar nathī jatī.}
\end{quote}

\begin{quote}
This lane is so narrow that a car does not go in.
\end{quote}

\begin{cousinweb}
\textbf{\gu{સાંકડું}} (\emph{sā̃kḍũ}) is a tadbhava of Sanskrit \emph{saṁkaṭa}, \textquotedblleft{}narrow, pressed together\textquotedblright{}; CDIAL lists the Gujarati word (CDIAL 12817).
\end{cousinweb}

\begin{grammarlens}[title={What you've built}]
One more word for getting around.
\end{grammarlens}

\subsection*{Your turn}
\begin{itemize}
\raggedright
  \item \emph{Say it:} \emph{sā̃kḍũ}
  \item \emph{Say it:} \emph{sā̃kḍũ}, once more
  \item \emph{Say it:} \emph{ā galī eṭlī sā̃kḍī che ke gāḍī andar nathī jatī.} — This lane is so narrow that a car does not go in.
  \item \emph{Recall:} say the Gujarati for a pavement, then the Gujarati for a trail, then say \emph{sā̃kḍũ} again
\end{itemize}

\begin{tcolorbox}[breakable,colback=teal!4,colframe=teal!35!black,title={Before you move on}]
What does \gu{સાંકડું} mean? (\textquotedblleft{}Narrow\textquotedblright{}.) Say it once more.
\end{tcolorbox}
\section[\texorpdfstring{\gu{ચઢાણ} (caḍhāṇ)}{ચઢાણ (caḍhāṇ)}]{\gu{ચઢાણ} (caḍhāṇ) — an uphill climb; an ascent}

\label{lesson:GU-C404-cadhan}



\begin{quote}
Before the new one: say the Gujarati for a trail, then the Gujarati for narrow.
\end{quote}

\subsection*{You'll want to know: \gu{ચઢાણ}}
\textbf{\gu{ચઢાણ}} — \emph{caḍhāṇ} — \textquotedblleft{}an uphill climb; an ascent\textquotedblright{}.

\begin{quote}
\gu{મંદિર} \gu{સુધીનું} \gu{ચઢાણ} \gu{ઘણું} \gu{અઘરું} \gu{છે}.
\end{quote}

\begin{quote}
\emph{mandir sudhīnũ caḍhāṇ ghaṇũ agharũ che.}
\end{quote}

\begin{quote}
The climb up to the temple is very hard.
\end{quote}

A neuter noun.

\begin{cousinweb}
\textbf{\gu{ચઢાણ}} (\emph{caḍhāṇ}) is built on \gu{ચઢવું}, \textquotedblleft{}to climb\textquotedblright{}, which goes back to \emph{*caḍhati}, \textquotedblleft{}rises\textquotedblright{}, a form CDIAL reconstructs (CDIAL 4578): the stretch of road you have to climb.
\end{cousinweb}

\begin{grammarlens}[title={What you've built}]
One more word for getting around.
\end{grammarlens}

\subsection*{Your turn}
\begin{itemize}
\raggedright
  \item \emph{Say it:} \emph{caḍhāṇ}
  \item \emph{Say it:} \emph{caḍhāṇ}, once more
  \item \emph{Say it:} \emph{mandir sudhīnũ caḍhāṇ ghaṇũ agharũ che.} — The climb up to the temple is very hard.
  \item \emph{Recall:} say the Gujarati for a trail, then the Gujarati for narrow, then say \emph{caḍhāṇ} again
\end{itemize}

\begin{tcolorbox}[breakable,colback=teal!4,colframe=teal!35!black,title={Before you move on}]
What does \gu{ચઢાણ} mean? (\textquotedblleft{}An uphill climb; an ascent\textquotedblright{}.) Say it once more.
\end{tcolorbox}
\section[\texorpdfstring{\gu{ટપવું} (ṭapvũ)}{ટપવું (ṭapvũ)}]{\gu{ટપવું} (ṭapvũ) — to jump over; to leap across}

\label{lesson:GU-C404-tapvu-leap}



\begin{quote}
Before the new one: say the Gujarati for narrow, then the Gujarati for an uphill climb.
\end{quote}

\subsection*{You'll want to know: \gu{ટપવું}}
\textbf{\gu{ટપવું}} — \emph{ṭapvũ} — \textquotedblleft{}to jump over; to leap across\textquotedblright{}.

\begin{quote}
\gu{છોકરો} \gu{ખાડો} \gu{ટપી} \gu{ગયો}.
\end{quote}

\begin{quote}
\emph{chokro khāḍo ṭapī gayo.}
\end{quote}

\begin{quote}
The boy jumped over the pothole.
\end{quote}

Past (he …): \gu{તે} \gu{ખાડો} \gu{ટપી} \gu{ગયો}, \emph{te khāḍo ṭapī gayo}.

\begin{cousinweb}
\textbf{\gu{ટપવું}} (\emph{ṭapvũ}) goes back to \emph{*tarpati}, \textquotedblleft{}jumps\textquotedblright{}, a form CDIAL reconstructs; CDIAL lists the Gujarati verb, \textquotedblleft{}to leap, to jump over\textquotedblright{} (CDIAL 5727).
\end{cousinweb}

\begin{grammarlens}[title={What you've built}]
One more word for getting around.
\end{grammarlens}

\subsection*{Your turn}
\begin{itemize}
\raggedright
  \item \emph{Say it:} \emph{ṭapvũ}
  \item \emph{Say it:} \emph{ṭapvũ}, once more
  \item \emph{Say it:} \emph{chokro khāḍo ṭapī gayo.} — The boy jumped over the pothole.
  \item \emph{Recall:} say the Gujarati for narrow, then the Gujarati for an uphill climb, then say \emph{ṭapvũ} again
\end{itemize}

\begin{tcolorbox}[breakable,colback=teal!4,colframe=teal!35!black,title={Before you move on}]
What does \gu{ટપવું} mean? (\textquotedblleft{}To jump over; to leap across\textquotedblright{}.) Say it once more.
\end{tcolorbox}
